% ===================================================================
%  GAMBAR No. 11 — Kutha lan Monumene. --- Kobongan.
%
%  Traduction, faite en 2026, en javanais d'aujourd'hui, sur l'ido de
%  texto/io. Regles de la colonne : voir 10-gambar-01.tex. Deux
%  scenes, deux numerotations. Ouverture de la TROISIEME SERIE :
%  « SERI KATELU ». Deux blocs marques « ¶ » par ordo.py.
%
%  C'EST LE TABLEAU QUI N'AVAIT AUCUN GROS PLAN DANS LES 42 COLONNES,
%  ET IL EN A MAINTENANT. La planche 11 numerote c1/c2 pour ses deux
%  scenes mais grave ses numeros NUS ; html.py cherchait « c1:7 » et
%  ne trouvait rien, dans toutes les langues a la fois et depuis
%  toujours. On ne s'en est apercu qu'en comptant les boutons tableau
%  par tableau, pendant l'ecriture de la colonne ourdoue. Le repli sur
%  le numero nu a rendu ses 89 gros plans a ce tableau — une colonne
%  complete est passee de 1 644 boutons a 1 733.
%
%  LA VILLE EST LE PREMIER TABLEAU DU LIVRET OU LE JAVANAIS EMPRUNTE
%  PLUS QU'IL NE GARDE, ET LA RAISON EST DANS LA DATE DES CHOSES. La
%  ville javanaise ancienne a ses mots — kutha la ville, alun-alun la
%  place, pasar, kraton, pendhapa, gapura, wates —, mais tout ce que
%  la planche montre est du XIXe siecle europeen : la douane, la
%  Bourse, la prefecture, la caisse d'epargne, le boulevard, le musee,
%  l'ecole normale, la poste, le tramway. Aucun de ces batiments
%  n'existait avant que le Hollandais les construise, et ils portent
%  donc son nom : pabean, bursa, kantor bupati, bank tabungan,
%  bulevar, museum, kantor pos, trem.
%
%  ET « pabean » (10) EST L'EXCEPTION QUI LE PROUVE. Le mot est
%  javanais et malais, de bea, la taxe : la douane existait aux ports
%  de Java bien avant les Hollandais, et elle est le seul de ces
%  batiments a garder un nom d'ici. Ce qui a une histoire ici garde
%  son mot ; ce qui est arrive tout fait garde le sien.
%
%  UN MOT INDONESIEN EST ICI LA SEULE FORME JUSTE, ET kolonoj.py A
%  DU L'APPRENDRE. L'hopital (39) se dit « rumah sakit » a Java comme
%  partout en Indonesie ; or « rumah » etait releve depuis le cablage
%  comme indonesianisme, le javanais disant omah. La regle avait
%  raison sur le MOT et tort sur le NOM : rumah sakit est le nom de
%  l'institution, et la seule autre forme possible — « griya sakit »
%  — est du KRAMA, que la regle de niveau releverait a son tour. La
%  regle a donc ete resserree, « rumah » sauf devant « sakit », et
%  verifiee en glissant expres « rumah gedhe » dans ce fichier : le
%  controle le releve, puis la ligne a ete otee. Un mot emprunte peut
%  etre la seule forme juste, et un controle doit savoir distinguer un
%  mot d'un nom.
%
%  L'INCENDIE, LUI, RETOURNE AU JAVANAIS D'UN SEUL COUP : geni le feu,
%  kobongan l'incendie, kukus la fumee (75), andha l'echelle (72),
%  ilat geni les flammes, umob. Une ville brule dans la langue de ceux
%  qui l'eteignent, meme quand ses batiments portent des noms
%  etrangers. La meme frontiere qu'au tableau 6 entre le pawon et le
%  kompor, et qu'au tableau 10 entre la voile et la vapeur : le mot
%  suit la CHOSE, jamais le prestige.
%
%  ET L'ORCHESTRE (80) FAIT ICI CE QUE LE TABLEAU 16 FERA POUR LE
%  THEATRE. Le javanais a « gamelan », le mot le plus javanais de tous,
%  et l'orchestre du theatre en flammes est une fosse d'opera avec
%  violon, flute, violoncelle, trombone et tambour. « orkestra » et
%  « biola » sont donc ecrits, et gamelan ne parait pas. C'est le
%  cinquieme refus de la colonne, apres gobak sodor, joged, gendhongan
%  et wayang — et le plus couteux, parce que c'est le mot que cette
%  langue aurait le plus de droit a poser.
% ===================================================================

%%K t11-apar-1 apar
\VUsaut{2.0mm}
\VUcentre{13.2pt}{90}{SERI KATELU}
\VUsaut{3.0mm}
\VUfilet{16mm}
\VUsaut{5.6mm}
\VUtitre{15.0pt}{60}{GAMBAR No. 11}
\VUsaut{1.4mm}
\VUpk{11.4pt}{40}{Kutha lan Monumene. --- Kobongan.}
\VUsaut{2.2mm}

%%K t11-tit-1 sub
\VUsaut{3.0mm}
\VUpk{10.2pt}{40}{Pinggir Kutha.}
\VUsaut{3.0mm}

%%K t11-c1-01-1 p
1. --- Aku manggon ing kutha gedhe kang dununge 100 kilometer saka
Samodra, lan kang sanajan mangkono dadi \VUgras{pelabuhan}
\textsuperscript{(1)} kang wigati. Kali kang mbatesi kutha iku ambane
400 meter lan mbentuk kelokan kang endah banget.

%%K t11-c1-02-1 p
2. --- Kabeh peranganing kutha kang ana ing \VUgras{dermaga}
\textsuperscript{(2)} katon sarwa laut lan sarwa pabrik, lan mbentuk
\VUgras{pinggir kutha} \textsuperscript{(3)} kang mung dienggoni wong laut
lan buruh. Wong-wong iku nyambut gawe ing \VUgras{pabrik}
\textsuperscript{(4)} lan ing \VUgras{pabrik gas} \textsuperscript{(5)},
kang \VUgras{cerobong}e \textsuperscript{(6)} kang dhuwur lan
\VUgras{gasometer}e katon saka adoh banget.

%%K t11-c1-03-1 p
3. --- Ing abad tengahan, kutha iku dibentengi, lan wong isih nemu
sawetara tilase tembok betenge, mligine \VUgras{gapura}
\textsuperscript{(8)} \VUgras{kastil beteng} \textsuperscript{(7)} lawas
kang diapit \VUgras{menara} \textsuperscript{(9)} loro kang saiki dianggo
\VUgras{pakunjaran.}

%%K t11-c1-04-1 p
4. --- Ing sakabehe \VUgras{kampung} iki, kang katon lawas banget, mung
siji bangunan kang rada anyar : yaiku \VUgras{pabean}
\textsuperscript{(10)}. Bangunan iku manggon ing antarane dermaga karo
\VUgras{dalan cilik} \textsuperscript{(11)} kang tumuju \VUgras{balai
kutha} \textsuperscript{(12)} lawas. Wong gampang ngenali monumen kang
keri iku merga \VUgras{menara lonceng} \textsuperscript{(13)} raseksa kang
nguwasani kutha kuna iku.

%%K t11-c1-05-1 p
5. --- Ing sisih sijine dalan, \VUgras{ngadeg} bangunan kang digawe
sagaya karo pabean : yaiku \VUgras{Bursa} \textsuperscript{(14)}. Salah
sijine ngarepane madhep alun-alun cilik panggonan \VUgras{kantor
bupati} \textsuperscript{(15)}. Nyatane, kutha kita, sanajan dudu
ibukutha gedhe, tetep dadi ibukutha kabupaten kang wigati.

%%K t11-tit-2 sub
\VUsaut{5.0mm}
\VUpk{10.2pt}{40}{Peranganing Kutha kang Paling Dhuwur.}
\VUsaut{3.4mm}

%%K t11-c1-06-1 p
6. --- Garnisune, kang cacahe akeh, manggon ing \VUgras{tangsi}
\textsuperscript{(16)} kang jembar lan sehat banget ; tangsi iku megar
tekan pager wesine \VUgras{taman kutha} \textsuperscript{(17)}.
\VUgras{Bulevar} \textsuperscript{(19)} kang tumuju taman iku liwat ing
mburine \VUgras{bank tabungan} \textsuperscript{(18)} lan tekan alun-alune
\VUgras{katedral} \textsuperscript{(20)}.

%%K t11-c1-07-1 p
7. --- Kabeh monumen iku megar tumumpang kaya amfiteater ing punthuk
cilik panggonan uga ana \VUgras{sekolah menengah} \textsuperscript{(21)}
kita kang jembar.

%%K t11-c1-07-2 p
Yen wong mudhun liwat dalan kang rada mendhak, kang nyambung karo Dalan
Gedhe, wong tekan \VUgras{Pengadilan} \textsuperscript{(22)}, kang ing
ngarepe megar \VUgras{taman umum} \textsuperscript{(26)} kang nyenengake.
\VUgras{Taman alun-alun} iku ora gedhe banget, nanging gawe ing
tengahing kutha oase ijo lan adhem. Mula taman iku rame banget
diparani wong kutha kang seneng lungguh ing sangisore tarube
\VUgras{warung kopi} \textsuperscript{(25)}, kanggo ngrungokake para
prajurit musik kang ing dina Kemis lan Minggu main ing
\VUgras{panggung bunder} \textsuperscript{(27)} kang diselehake ing
antarane palemahan suket lan gerumbulan.

%%K t11-c1-08-1 p
8. --- Ing salah sijine palemahan suket iku ngadeg \VUgras{reca}
\textsuperscript{(28)} prunggu, monumen kang diadegake kanggo pangeling
salah sijine wong sakutha kita, kang wis gawe jasa gedhe marang kutha
kelairane.

%%K t11-tit-3 sub
\VUsaut{4.0mm}
\VUpk{10.2pt}{40}{Piranti Angkutan.}
\VUsaut{3.4mm}

%%K t11-c1-09-1 p
9. --- Nganti saiki kutha kita durung dipadhangi lampu listrik ; kutha
kita mung duwe \VUgras{lampu gas} \textsuperscript{(29)}. Nanging
\VUgras{koran-koran ing kene} nganakake kampanye supaya sistem padhang
iku dibecikake, \VUgras{kaya} wong wis nyampurnakake \VUgras{sistem
tarikan trem.} \VUgras{Kreta} lawas, kang \VUgras{ditarik}
\VUgras{jaran,} sacara praktis wis diganti \VUgras{trem listrik}
\textsuperscript{(31)} kang cepet nggawa para lelungan saka pucuking
\VUgras{kutha menyang pucuk sijine,} kanthi rega kang \VUgras{murah}
banget.

%%K t11-c1-10-1 p
10. --- Cedhak Pengadilan ana \VUgras{Bank} \textsuperscript{(32)}
panggonan wong ngedhunake surat dagang. \VUgras{Juru tagih bank}
\textsuperscript{(33)} lunga nagih utang menyang omahe wong.

%%K t11-tit-4 sub
\VUsaut{4.0mm}
\VUpk{10.2pt}{40}{Seni lan Kawruh.}
\VUsaut{3.4mm}

%%K t11-c1-11-1 p
11. --- Bangunan endah kang ngadeg ing cedhake yaiku \VUgras{museum.}
Museum iku bebarengan isi \VUgras{perpustakaan} \textsuperscript{(34)}
kutha, \VUgras{klumpukan ngelmu alam} \textsuperscript{(35)} kang endah,
lan \VUgras{museum lukisan} \textsuperscript{(36)} kang apik, kang dadi
\VUgras{pameran} lestari kang endah banget.

%%K t11-c1-11-2 p
Museum iku dibukak kanggo umum kaping pindho saben minggu, nanging wong
manca kena nyowani saben dina kanthi menehi \VUgras{juru jaga}ne
\textsuperscript{(37)} sethithik. \VUgras{Tukang lukis}
\textsuperscript{(38)} padha teka mrana nyambut gawe. Wong weruh
wong-wong iku \VUgras{ing ngarepe} sesangganing lukisan, karo
\VUgras{palet} ing tangan, nyalin lukisan kang misuwur.

%%K t11-c1-12-1 p
12. --- Ing panggonan kang dhuwur, ing mburine museum, wong wiwit weruh
\VUgras{kubah} \textsuperscript{(40)} \VUgras{rumah sakit}
\textsuperscript{(39)} lan bangunane \VUgras{sekolah guru}
\textsuperscript{(41)} panggonan para guru sekolah desa kita dilatih.

Ing alun-alune Teater ana \VUgras{kantor pos} \textsuperscript{(43)} kang
manggon ing \VUgras{omah pribadi} \textsuperscript{(42)}.

%%K t11-tit-5 sub
\VUsaut{5.0mm}
\VUpk{11.4pt}{40}{Kobongan.}
\VUsaut{3.0mm}

%%K t11-tit-6 sub
\VUsaut{3.4mm}
\VUpk{10.2pt}{40}{Bengok Kobongan.}
\VUsaut{3.4mm}

%%K t11-c2-01-1 p
1. --- Kabar kang medeni sumebar ing sakutha : lagi wae ana kobongan
ing teater! Pas nalika aku tekan \VUgras{alun-alun}
\textsuperscript{(96)}, wong akeh wis nglumpuk ing \VUgras{trotoar}
\textsuperscript{(46)} lan ing \VUgras{dalan gedhe} \textsuperscript{(47)}.
\VUgras{Tukang pos} \textsuperscript{(44)} lan \VUgras{tukang layang}
\textsuperscript{(45)} padha metu saka kantor pos. \VUgras{Sopir trem}
\textsuperscript{(48)} kepeksa mandhegake kretane supaya \VUgras{kreta
ambulans} \textsuperscript{(50)} kutha bisa liwat, kang teka karo
\VUgras{ahli bedah} \textsuperscript{(52)} lan \VUgras{juru rawat}
\textsuperscript{(51)} kanggo ngrawat \VUgras{wong tatu}
\textsuperscript{(53)} kang digawa nganggo \VUgras{tandhu}
\textsuperscript{(54)} menyang cedhake \VUgras{kios koran}
\textsuperscript{(55)}.

%%K t11-c2-02-1 p
2. --- Kompi infanteri, kang lagi bali saka latihan, mandheg kanggo
njaga tata tentrem bareng karo \VUgras{polisi kutha kang nunggang
jaran} \textsuperscript{(61)}. \VUgras{Para kang nunggang jaran,} kang
jarane padha njranthal, luwih becik tinimbang \VUgras{prajurit}
\textsuperscript{(57)} utawa \VUgras{polisi} \textsuperscript{(63)}
anggone mundurake \VUgras{wong kang padha kepengin weruh}
\textsuperscript{(56)}. Kejaba iku polisi padha repot banget. Salah
sijine nuduhake marang \VUgras{polisi kutha} \textsuperscript{(64)}
menyang ngendi wong kudu nggawa \VUgras{korban} \textsuperscript{(65)}
liyane, yaiku tukang pompa geni kang kendel kang tiba saka andha.

%%K t11-c2-03-1 p
3. --- \VUgras{Opsir} \textsuperscript{(66)} netepake panggonan
\VUgras{pompa uap} \textsuperscript{(67)} kang lagi teka kudu diselehake.
Karo ngenteni mesin kang kuwat iku, dheweke enggal-enggal masang
\VUgras{pompa tangan} \textsuperscript{(68)} kang \VUgras{selang}e
\textsuperscript{(69)} kang dawa nyemprotake banyu kaya grojogan menyang
geni.

Wong nem tukang pompa ngobahake pompa iku, dene kancane, kang nyekel
\VUgras{muncrat}e \textsuperscript{(70)}, ngarahake \VUgras{semprotan}
\textsuperscript{(71)} menyang tengahing kobongan.

%%K t11-c2-04-1 p
4. --- Ing sisih sijine teater, wong wis nyanggongake \VUgras{andha}
\textsuperscript{(72)} gedhe. Andha iku mbisakake para tukang pompa
nyedhaki ilat geni kang muncrat ing antarane ubengan \VUgras{kukus}
\textsuperscript{(75)} ireng.

%%K t11-tit-7 sub
\VUsaut{5.0mm}
\VUpk{10.2pt}{40}{Lelabuhan kang Kendel.}
\VUsaut{3.4mm}

%%K t11-c2-05-1 p
5. --- \VUgras{Kobongan} \textsuperscript{(74)} iku lair ing ruwang
ngarepe \VUgras{teater} \textsuperscript{(73)}, nalika wong lagi gladhi
\VUgras{opera} kang \VUgras{pagelaran} kapisane bakal dianakake sesuk.
Begjane mung sethithik \VUgras{wong nonton} \textsuperscript{(76)} kang
ana ing ruwang tontonan. Wong-wong kang mesakake iku ngungsi ing
\VUgras{balkon} \textsuperscript{(77)}, panggonan \VUgras{tukang pompa
geni} \textsuperscript{(86)} kang kendel teka nulungi nganggo andha lan
tampar.

%%K t11-c2-06-1 p
6. --- Ing sangisore \VUgras{peristil} \textsuperscript{(78)}, panggonan
\VUgras{plakat} \textsuperscript{(79)} aweh kabar bab pagelaran, wong
weruh \VUgras{para pemain musik} \textsuperscript{(81)}
\VUgras{orkestra} \textsuperscript{(80)} padha mlayu karo nggawa
alat-alate : \VUgras{biola} \textsuperscript{(81)}, \VUgras{suling}
\textsuperscript{(82)}, \VUgras{selo} \textsuperscript{(83)},
\VUgras{trombon} \textsuperscript{(84)}, \VUgras{tambur}
\textsuperscript{(85)}, lan sapanunggalane. Wong ngudi nylametake
peranganing \VUgras{perabot} \textsuperscript{(87)}.

%%K t11-c2-07-1 p
7. --- \VUgras{Aktor lanang} \textsuperscript{(89)} lan \VUgras{aktor
wadon} \textsuperscript{(88)}, \VUgras{penyanyi lanang} lan
\VUgras{penyanyi wadon} kepeksa mlayu karo \VUgras{sandhangan}
\textsuperscript{(90)} panggunge. Sang tenor nerangake marang
\VUgras{wartawan} \textsuperscript{(92)} kepriye kedadeyane.
\VUgras{Juru warta} sanalika mlayu teka wiwit wektu kapisan bareng
karo para pejabat : \VUgras{bupati} \textsuperscript{(91)},
\VUgras{wali kutha} \textsuperscript{(93)}, \VUgras{komisaris polisi}
\textsuperscript{(94)} lan \VUgras{jaksa} \textsuperscript{(95)}, kang
bakal nliti kanggo mutusake sapa kang tanggung jawab.

\VUsaut{6.0mm}
\VUfilet{22mm}
